\documentclass[11pt,a4paper]{article}
\usepackage[margin=27mm]{geometry}
\usepackage[T1]{fontenc}
\usepackage{lmodern}
\usepackage{amsmath,amssymb,amsthm}
\usepackage{booktabs,array}
\usepackage[hidelinks]{hyperref}
\hypersetup{pdftitle={Quadratic maps satisfying f composed with f = f + id: polar Lie structure and sharp dimension bounds},pdfauthor={}}
\renewcommand{\labelenumi}{\textup{(\roman{enumi})}}
\setlength{\emergencystretch}{2em}
\newtheorem{theorem}{Theorem}[section]
\newtheorem{proposition}[theorem]{Proposition}
\newtheorem{lemma}[theorem]{Lemma}
\newtheorem{corollary}[theorem]{Corollary}
\theoremstyle{definition}
\newtheorem{definition}[theorem]{Definition}
\newtheorem{example}[theorem]{Example}
\theoremstyle{remark}
\newtheorem{remark}[theorem]{Remark}
\newcommand{\F}{\mathbb F}
\newcommand{\kfour}{\F_4}
\newcommand{\id}{\operatorname{id}}
\newcommand{\im}{\operatorname{im}}
\newcommand{\End}{\operatorname{End}}
\newcommand{\Hom}{\operatorname{Hom}}
\newcommand{\GL}{\operatorname{GL}}
\newcommand{\Aut}{\operatorname{Aut}}
\newcommand{\ad}{\operatorname{ad}}
\newcommand{\Span}{\operatorname{span}}
\newcommand{\Ann}{\operatorname{Ann}}
\newcommand{\rank}{\operatorname{rank}}
\newcommand{\rad}{\operatorname{rad}}
\newcommand{\cl}{\operatorname{cl}}
\newcommand{\m}{\mathfrak m}
\DeclareMathOperator{\supp}{supp}
\title{Quadratic maps satisfying \(f\circ f=f+\id\):\\
polar Lie structure and sharp dimension bounds}
\author{}
\date{}
\begin{document}
\maketitle
\begin{abstract}
The equation \(f\circ f=f+\id\) couples three-cycles to vector addition:
each nonzero orbit consists of the three nonzero points of a binary plane.
We study its solutions of binary algebraic degree at most two through their
polar Lie algebras. The equation forces both the Jacobi identity and an
additional degree-four identity, and every solvable polar algebra is nilpotent.
In the metabelian case we obtain a complete normal form and an equivalence,
including isomorphisms and automorphism groups, with metabelian
\(\F_4\)-Lie algebras whose quotient actions on the derived ideal have square zero.
The correspondence gives the exact minimum binary dimension at every polar
nilpotency class \(c\ge2\): \(6,8,12\) in classes two, three, and four,
and \(2^c+2c-6\) for \(c\ge5\).
We characterize when the original polar bracket admits an \(\F_4\)-bilinear
structure extending its intrinsic scalar actions. In this compatible subclass
the minimum is \(2^c+2c-2\) for every \(c\ge2\), with at least \(c-2\)
inequivalent extremizers for every \(c\ge4\).
Two closed classifications illustrate the structure theory: all
eight-dimensional solutions form three linear-conjugacy classes, and the
solvable solutions with four-dimensional binary abelianization form an
exhaustive chain of seven classes ending in dimension sixteen. Exact central-extension
certificates make the finite rank calculations in the latter classification
independently reproducible.
\end{abstract}
\noindent\textbf{Keywords.}
Quadratic map; complete permutation; polar Lie algebra; metabelian Lie algebra;
restricted Lie algebra; nilpotency class; finite field.

\section{Introduction}
\label{sec:introduction}

Let \(V\) be a finite-dimensional vector space over \(\F_2\). We consider
the functional equation
\begin{equation}\label{eq:functional}
 f(f(x))=f(x)+x\qquad(x\in V)
\end{equation}
under the assumption that \(f\) has binary algebraic degree at most two.
Every nonzero orbit is the triple \(\{x,f(x),x+f(x)\}\), the nonzero
points of a binary plane. The equation therefore ties a local cycle
condition to the additive geometry of the whole space. We ask how a
quadratic degree bound organizes these orbits into an algebraic structure
that can be classified and measured.

The linear case identifies the relevant scalar field. A linear solution
has minimal polynomial \(X^2+X+1\) on every nonzero space, so it defines
an intrinsic \(\F_4\)-structure in which \(f\) is multiplication by
\(\omega\), where \(\omega^2+\omega=1\). Up to binary-linear conjugacy,
there is one such solution in each positive even dimension.
Degree two is the first level at which the polar form can be nonzero.
It introduces interactions between the orbit planes, and these
interactions are the subject of the paper.

The first structural object is the polar form
\[
 [x,y]=f(x+y)+f(x)+f(y).
\]
For a general zero-constant quadratic map this form is alternating and
bilinear. Here the functional equation also forces Jacobi, making it a
Lie bracket. The derived space, the center, and the lower central series
therefore become intrinsic invariants of the map. The equation imposes
further rigidity: solvability implies nilpotency, and fourth polarization
gives a degree-four polynomial identity. This polar structure is the
bridge from the permutation equation to both classification and sharp
dimension bounds.

\subsection{The main correspondence}

Write \(W=[V,V]\) for the binary span of all polar values.
The principal classification theorem concerns the metabelian case
\([W,W]=0\). The maps induced by \(f\) on \(W\) and \(V/W\) are linear
operators satisfying \(T^2+T=\id\); they therefore define intrinsic
\(\F_4\)-vector space structures. Fix
\(\omega\in\F_4\) with \(\omega^2+\omega=1\).

In suitable binary-linear coordinates, every metabelian solution has the form
\begin{equation}\label{eq:intro-normal}
 f(x,z)=\bigl(\omega x,\,
       \omega z+P(x)z+Q(x)+P(x)Q(x)\bigr),
\end{equation}
where
\[
 P(x)=\sum_{i=1}^d x_iP_i,
 \qquad Q(x)=\sum_{i<j}x_ix_ja_{ij}.
\]
The \(\F_4\)-linear operators \(P_i\) commute, satisfy \(P_i^2=0\), and obey
\[
 P_i a_{jk}+P_j a_{ik}+P_k a_{ij}=0.
\]
The coefficients \(a_{ij}\) generate \(W\) under the operators \(P_i\).
Conversely, all such data give solutions.

The key necessity statement is that binary degree two forces the
square-zero actions on the intrinsic polar image \(W\). Thus the normal
form describes every metabelian solution as well as constructing them.
Moreover, the correspondence includes morphisms. Binary-linear conjugacy
of these maps is equivalent to \(\F_4\)-Lie isomorphism of the associated
metabelian algebras with \((\ad u)^2|_{W}=0\).
In particular it identifies the linear centralizer of \(f\) with the
automorphism group of the associated \(\F_4\)-Lie algebra
(Theorem~\ref{thm:correspondence}).

\subsection{Sharp dimension costs}

The normal form gives an exact measure of the dimension required for a
given polar Lie nilpotency class. Let \(\mu(c)\) be the minimum binary
dimension among metabelian solutions of class \(c\ge2\). We prove
\begin{equation}\label{eq:intro-minimum}
 \mu(2)=6,\qquad\mu(3)=8,\qquad\mu(4)=12,\qquad
 \mu(c)=2^c+2c-6\quad(c\ge5).
\end{equation}
The proof uses a universal commutator module over a square-zero polynomial
ring. Its multigrading determines the possible terminal commutators, and a
dual contraction argument determines the smallest quotient that retains
one of them. Both the lower bound and its attainment are uniform in \(c\)
(Theorem~\ref{thm:unrestricted-minimum}).

The \(\F_4\)-Lie algebra in this correspondence is auxiliary: the original
binary polar bracket need not itself be \(\F_4\)-bilinear. We give an exact
linear criterion for when it admits such a structure with the prescribed
intrinsic scalar actions on \(W\) and \(V/W\).
This condition is equivalent to the existence of a restricted square on
the auxiliary algebra that vanishes on \(W\) and takes values in \(W\)
(Theorem~\ref{thm:compatibility}).

Let \(\nu(c)\) denote the minimum dimension in this compatible subclass.
We obtain the single formula
\begin{equation}\label{eq:intro-compatible-minimum}
 \nu(c)=2^c+2c-2\qquad(c\ge2).
\end{equation}
One part of its lower bound counts independent preimages of the nonempty
faces of a square-zero operator cube. This gives
\(\dim_{\F_4}W\ge2^{c-1}-1\), independently of the unrestricted extremal
argument. The other part proves \(\dim_{\F_4}(V/W)\ge c\).
The elementary augmentation-module construction attains both bounds
(Theorem~\ref{thm:compatible-minimum}).

\begin{table}[ht]
\centering
\caption{Minimum binary dimensions in the two metabelian subclasses.}
\label{tab:minima}
\begin{tabular}{@{}lrrrrrr@{}}
\toprule
Polar class \(c\)&2&3&4&5&6&7\\
\midrule
\(\mu(c)\)&6&8&12&36&70&136\\
\(\nu(c)\)&6&12&22&40&74&140\\
\bottomrule
\end{tabular}
\end{table}

Compatibility therefore has an exact additional cost of four binary
dimensions in every class \(c\ge5\). Equality still allows growing
structural diversity: for each \(c\ge4\) we construct at least \(c-2\)
pairwise inequivalent compatible extremizers
(Theorem~\ref{thm:extremizers}).

\subsection{Closed classification ranges}

Two applications illustrate the reach and the boundaries of the theory.
There are exactly three binary-linear conjugacy classes in dimension eight
(Theorem~\ref{thm:dimension-eight}). The proof combines the correspondence
with the classical small-dimensional simple Lie algebra lists and short
explicit obstructions to the polar identities.

There is also a closed classification in an all-dimensional stratum:
if the polar algebra is solvable and
\(\dim_{\F_2}(V/W)=4\), exactly one class occurs in each of dimensions
\(4,6,8,10,12,14,16\), and none occurs thereafter
(Theorem~\ref{thm:seven-chain}). Its final two members are nonmetabelian.
Completeness is proved by a lossless central descent, exact extension-space
calculations, and explicit projective orbit actions. The finite rank portion
is supplied with independently replayable certificates.

\subsection{Context and related work}

There is an operator viewpoint on \eqref{eq:functional}. On complete
mappings, \(\mathcal F(g)=(\id+g)^{-1}\) satisfies
\(\mathcal F^3(g)=g\), and its fixed objects are exactly the solutions
of \eqref{eq:functional}. The arithmetic study of norm fibers and affine
companion groups in~\cite{Li2026} also describes these fixed objects by
oriented binary line spreads. The present paper studies their quadratic
structure through the polar Lie algebra.

Equation~\eqref{eq:functional} implies that both \(f\) and \(f+\id\) are
permutations. Triple-cycle and regular complete permutation polynomials
have been studied by Liu--Chen--Xu--Sun~\cite{Liu2019},
Xu--Zeng--Zhang~\cite{Xu2022}, and Lu--Wu--Wang--Cao~\cite{Lu2023}.
Quadratic complete mappings also occur in the constructions of
Samardjiska--Gligoroski~\cite{Samardjiska2017}.
Here the orbit-sum identity, binary degree, and binary-linear conjugacy
are considered together. In particular, simultaneous three-cycle behavior
of \(f\) and \(f+\id\) does not imply \eqref{eq:functional};
Example~\ref{ex:regular-cpp} gives a quadratic instance from the
affine-switching construction of~\cite{Xu2022}.

There is a parallel quasigroup interpretation. The operation
\(x*y=x+f(x+y)\) is a quaternary Mendelsohn quasigroup.
Zhang's thesis~\cite[Chapter~3]{Zhang1998} develops central representations
and recursions in that setting. Its Lemma~10.3(1) already contains the
central six-dimensional example used below. Throughout this paper,
equivalence means binary-linear conjugacy, and nilpotency class means
the class of the polar Lie algebra. These specifications retain the
addition and the degree filtration that enter the necessity proofs.

The commutator-module presentation is the classical Koszul presentation
of the free metabelian derived module; see
Papadima--Suciu~\cite[Theorem~6.2]{PapadimaSuciu2004} and the
characteristic-two module methods of Kov\'acs--St\"ohr~\cite{KovacsStohr2013}.
Restricted squares and quadratic maps have a broader established
connection; see, for example, Pakianathan--Yal\c{c}{\i}n~\cite{Pakianathan2007}.
We use these tools to prove the degree-forced correspondence and the exact
extremal formulas for \eqref{eq:functional}. The four-variable polar identity
is the classical degree-four standard polynomial in Lie form
\cite[Theorem~9]{Drensky2021}; its forced occurrence here provides the
obstructions used in the low-dimensional application.

\subsection{Conventions}

We use \(f^j\) only for composition powers. All dimensions without a
specified field are binary dimensions. A map is called quadratic when its
binary algebraic degree is at most two, so the linear cases are included.
Two solutions are equivalent if \(Sf=f'S\) for a binary-linear isomorphism
\(S\). Lie brackets are alternating; in characteristic two they are
symmetric as bilinear maps. We put
\(\gamma_1(L)=L\), \(\gamma_{i+1}(L)=[L,\gamma_i(L)]\), and
\(L^{(0)}=L\), \(L^{(i+1)}=[L^{(i)},L^{(i)}]\).
The class of a nonzero abelian algebra is one.

\medskip
\noindent\textbf{Notation at a glance.}
The \(\F_4\)-dimensions below are used in the metabelian correspondence.
\begin{center}
\small
\renewcommand{\arraystretch}{1.12}
\begin{tabular}{@{}p{.25\textwidth}p{.71\textwidth}@{}}
\toprule
Symbol & Meaning\\
\midrule
\(L=(V,[\ ,\ ])\) & The binary polar Lie algebra of \(f\).\\
\(W=L',\ U=V/W\) & The intrinsic derived ideal and abelianization.\\
\(d,\ w\) & \(\dim_{\F_4}U\) and \(\dim_{\F_4}W\).\\
\(\omega\) & A root of \(X^2+X+1\) in \(\F_4\).\\
\(\m\) & The augmentation ideal \((t_1,\ldots,t_d)\) of \(R_d\),
or its image under the operator action.\\
\(\gamma_i(L),\ L^{(i)}\) & The lower central and derived series.\\
\(\mu(c),\ \nu(c)\) & Minimum binary dimensions at class \(c\):
metabelian and compatible metabelian, respectively.\\
\(\mathcal A_k\) & Filtered action spaces used to distinguish extremizers
in Theorem~\ref{thm:extremizers}.\\
\bottomrule
\end{tabular}
\end{center}

\section{The polar Lie algebra}
\label{sec:polar}

\begin{proposition}\label{prop:cycles}
Every solution of \eqref{eq:functional} satisfies
\[
 f^3=\id,\qquad f(0)=0,\qquad
 f(x+f(x))=x.
\]
Its only fixed point is zero. Every nonzero orbit spans a binary plane,
and these planes partition \(V\setminus\{0\}\). In particular \(\dim V\)
is even.
\end{proposition}
\begin{proof}
Evaluate \eqref{eq:functional} at \(f(x)\). It gives
\(f^3(x)=f^2(x)+f(x)=x\), so \(f\) is bijective.
If \(a=f(0)\), then \(f(a)=a\); applying \(f^3\) at zero gives \(a=0\).
A fixed point in \eqref{eq:functional} must be zero.
The remaining identity follows from
\(x+f(x)=f^2(x)\). The nonzero orbits have length three and sum zero,
and \(3\mid 2^{\dim V}-1\) implies the parity assertion.
\end{proof}

\begin{theorem}\label{thm:polar}
Let \(f\) be a quadratic solution. Its polar form
\([x,y]=f(x+y)+f(x)+f(y)\) is a Lie bracket. It satisfies
\begin{align}
 [x,f(x)]&=0,\label{eq:commuting}\\
 f([x,y])&=[f(x),f(y)]+[x,y]
                  +[f(x)+f(y),[x,y]],\label{eq:bracket-image}\\
 [[a,b],[c,d]]+[[a,c],[b,d]]+[[a,d],[b,c]]&=0.
 \label{eq:four-variable}
\end{align}
\end{theorem}
\begin{proof}
The polar is alternating and bilinear because \(f\) is quadratic and
\(f(0)=0\). Expand \(f(x+f(x))=x\) by polarization. The terms
\(f(x)+f^2(x)\) sum to \(x\), leaving \eqref{eq:commuting}.
Take the third additive difference at zero of \(x\mapsto[x,f(x)]\).
The terms involving individual values of \(f\) cancel, leaving
\[
 [x,[y,z]]+[y,[z,x]]+[z,[x,y]]=0.
\]
This is Jacobi.

Next polarize \(f^2=f+\id\). Since
\(f(x+y)=f(x)+f(y)+[x,y]\), the expansion of the left side gives
\eqref{eq:bracket-image}.

For \eqref{eq:four-variable}, introduce Boolean variables
\(t_1,\ldots,t_4\) and vectors \(x_1,\ldots,x_4\). The exact expansion is
\[
 f\Bigl(\sum_i t_ix_i\Bigr)
   =\sum_i t_if(x_i)+\sum_{i<j}t_it_j[x_i,x_j].
\]
In applying \(f\) again, the coefficient of \(t_1t_2t_3t_4\) can only
come from a polar cross term between two disjoint pair summands.
There are three such pairings, giving the left side of
\eqref{eq:four-variable}. The composition has degree at most two,
so this coefficient vanishes. All coefficient comparisons take place in
the Boolean function algebra, with \(t_i^2=t_i\).
\end{proof}

The two polarization steps have distinct roles. Binary degree two gives
an alternating bilinear form; the commuting identity
\eqref{eq:commuting} forces its Jacobi identity. The fourth polarization
then supplies \eqref{eq:four-variable}, an additional restriction on the
resulting Lie algebra.
The polynomial in \eqref{eq:four-variable} is the associative standard
polynomial of degree four expressed as a Lie element in characteristic
two~\cite[Theorem~9]{Drensky2021}. Its forced occurrence here gives the
explicit exclusions in Section~\ref{sec:classification}.

\begin{proposition}\label{prop:invariant-series}
All terms of the derived, lower central, and upper central series of the
polar Lie algebra \(L=(V,[\ ,\ ])\) are \(f\)-invariant.
The center is \(f\)-invariant and \(f\) is linear on it.
Every nonzero factor of any of these series has even binary dimension.
\end{proposition}
\begin{proof}
Let \(I\) be an \(f\)-invariant ideal. In
\eqref{eq:bracket-image}, take \(x\in I\). Each term on the right belongs
to \([I,L]\), which is an ideal. Quadratic expansion extends invariance
from bracket generators to their span: the extra polar terms also lie in
that ideal. The same argument with \(x,y\in I\) proves invariance of
\([I,I]\), using that it is an ideal of \(L\).
Induction starts with \(I=L\).

Polarization of \eqref{eq:commuting} gives
\begin{equation}\label{eq:commuting-polarized}
 [x,f(y)]+[y,f(x)]+[x,[x,y]]+[y,[x,y]]=0.
\end{equation}
If \(x\in Z(L)\), this says \([y,f(x)]=0\) for every \(y\).
Thus \(Z(L)\) is invariant; its polar is zero, so \(f\) is linear there.
Pass successively to quotients by invariant central ideals to obtain
the upper central assertion. Every series factor inherits a solution of
\eqref{eq:functional}, so Proposition~\ref{prop:cycles} proves its parity.
\end{proof}

\begin{lemma}[Ideal descent]\label{lem:ideal-descent}
For every Lie ideal \(I\) of \(L\),
\[
 \{f(x)+I:x\in I\}\subseteq Z(L/I).
\]
Every perfect ideal is \(f\)-invariant. Every ideal \(I\) for which
\(Z(L/I)=0\) is also \(f\)-invariant.
\end{lemma}
\begin{proof}
For \(x\in I\), all terms in \eqref{eq:commuting-polarized} except
\([y,f(x)]\) belong to \(I\). This proves the inclusion.
The map \(\phi:I\to L/I\), \(x\mapsto f(x)+I\), is binary-linear
because its polar takes values in \(I\).
Equation~\eqref{eq:bracket-image} gives
\(\phi([x,y])=[\phi(x),\phi(y)]\) for \(x,y\in I\).
Its image is central, so \(\phi([I,I])=0\). This proves the perfect-ideal
assertion. The centerless-quotient assertion follows immediately from
the inclusion. Bijectivity and finiteness turn each invariant inclusion
into equality.
\end{proof}

\begin{corollary}\label{cor:simple-reduction}
If a quadratic solution has a nonsolvable polar algebra, then there is a
quadratic solution whose polar algebra is nonabelian simple, of even
binary dimension no larger than the original dimension.
\end{corollary}
\begin{proof}
The stabilized derived ideal \(P\) is nonzero and perfect, hence invariant.
Choose a maximal proper Lie ideal \(M\) of \(P\). The quotient \(P/M\)
is nonabelian simple and centerless. Apply Lemma~\ref{lem:ideal-descent}
inside \(P\) to see that \(M\) is invariant. Restriction and quotient
therefore give the required quadratic solution.
\end{proof}

\begin{lemma}\label{lem:two-generators}
A nonabelian nilpotent polar Lie algebra has binary abelianization
dimension at least four.
\end{lemma}
\begin{proof}
The abelianization is a nonzero even-dimensional space with the linear
action induced by \(f\). If its dimension were two, any nonzero coset
\(x+L'\) and its image \(f(x)+L'\) would form a binary basis.
The subalgebra \(H\) generated by \(x,f(x)\) is abelian by
\eqref{eq:commuting}. Moreover \(L=H+L'\).
Inductively, \(L=H+\gamma_j(L)\) implies
\(L=H+\gamma_{j+1}(L)\), because
\([H+\gamma_j,H+\gamma_j]\subseteq H+\gamma_{j+1}\).
Nilpotency then gives \(L=H\), a contradiction.
\end{proof}

\begin{example}\label{ex:regular-cpp}
Let \(\F_{16}\) contain \(\F_4\), let \(\eta\in\F_4\setminus\F_2\),
and write \(\chi(x)=1\) off \(\F_4\) and \(\chi(x)=0\) on \(\F_4\).
The affine-switching construction of
\cite[Corollary~6]{Xu2022} includes
\[
 G(x)=\eta x+\chi(x)
      =\eta x+\operatorname{Tr}_{\F_{16}/\F_4}(x)^3.
\]
The trace is binary-linear, and its cube is binary-quadratic.
Both \(\F_4\) and its complement are invariant under multiplication
by \(\eta\) and under translation by one. Hence \(G\) is a permutation
with \(G^3=\id\) and only zero fixed. The same holds for \(G+\id\).
Nevertheless
\[
 G^2(x)+G(x)+x=\eta\chi(x),
\]
so \(G\) does not satisfy \eqref{eq:functional}.
This separates the equation studied here from simultaneous triple-cycle
complete permutation behavior even within binary degree two.
\end{example}

\section{A complete metabelian correspondence}
\label{sec:metabelian}

Assume in this section that the polar Lie algebra is metabelian.
Then \(W=L'\) is abelian, \(f|_W\) is linear, and the induced map on
\(U=V/W\) is linear. Both operators satisfy \(T^2+T=\id\).
We give \(U\) and \(W\) the \(\F_4\)-structures in which these operators
are multiplication by \(\omega\).

\begin{lemma}[The operator pencil]\label{lem:operator-pencil}
In any binary-linear splitting of \(V\to U\), write
\begin{equation}\label{eq:preliminary-form}
 f(x,z)=(\omega x,\omega z+L_xz+q(x)).
\end{equation}
The binary-linear pencil \(x\mapsto L_x\) satisfies
\[
 L_x\omega=\omega L_x,\quad L_x^3=0,\quad
 L_{\omega x}=\omega L_x+L_x^2,\quad
 L_xL_y=L_yL_x,\quad L_x^2L_y=0.
\]
\end{lemma}
\begin{proof}
Put \(J=\omega\id_W\) and \(M_x=J+L_x\).
Both \(M_x\) and \(\id+M_x\) are invertible: they are the restrictions
on coset differences of the permutations \(f\) and \(f+\id=f^2\).
Substitution in \eqref{eq:functional} gives
\[
 M_{\omega x}M_x=\id+M_x.
\]
For \(A=M_x\), the three matrices along the \(\omega\)-orbit are
\(A,\id+A^{-1},(\id+A)^{-1}\). Their sum is \(J\), because the
three \(L\)-terms sum to zero. Thus \(J\) is a rational function of
\(A\), and in particular commutes with it. Clearing the denominators gives
\[
 JA(\id+A)=A^3+A+\id.
\]
Since \(J^2+J=\id\) and \(J^3=\id\), this is
\((A+J)^3=0\). It follows that \(L_x\) commutes with \(J\) and is
cube zero. Expanding \((J+L_x)^{-1}\) now gives
\(L_{\omega x}=JL_x+L_x^2\).
The additivity of the left side forces the \(L_x\)'s to commute.

For \(P=L_x,Q=L_y\), the identity \((P+Q)^3=0\) yields
\(P^2Q=PQ^2=:T\). Apply the same identity to
\(P'=JP+P^2=L_{\omega x}\) and \(Q\). It gives
\((J^2+J)T=P^2Q^2\), hence \(T=P^2Q^2=PT\).
As \(P\) is nilpotent, this implies \(T=0\).
\end{proof}

\begin{theorem}[Normal form and equivalence]\label{thm:correspondence}
Every metabelian quadratic solution is binary-linearly conjugate to
\begin{equation}\label{eq:normal-form}
 f(x,z)=\bigl(\omega x,\omega z+P(x)z+Q(x)+P(x)Q(x)\bigr)
       \quad\text{on }\F_4^d\oplus W,
\end{equation}
where
\begin{align}
 P(x)&=\sum_i x_iP_i,& Q(x)&=\sum_{i<j}x_ix_ja_{ij},\label{eq:PQ}\\
 P_iP_j&=P_jP_i,&P_i^2&=0,\label{eq:square-zero}\\
 P_i a_{jk}+P_j a_{ik}+P_k a_{ij}&=0\quad(i<j<k),
 \label{eq:triangle}\\
 W&=\sum_iP_iW+\Span_{\F_4}\{a_{ij}:i<j\}.
 \label{eq:intrinsic-generation}
\end{align}
Conversely these conditions define a quadratic solution with intrinsic
polar image \(W\).

The binary-linear conjugacy groupoid of these maps is equivalent to the
isomorphism groupoid of finite-dimensional metabelian \(\F_4\)-Lie
algebras \(\mathfrak l\) satisfying
\begin{equation}\label{eq:Lie-action-square}
 (\ad u)^2|_{\mathfrak l'}=0\qquad(u\in\mathfrak l).
\end{equation}
For corresponding objects,
\(C_{\GL(V,2)}(f)\cong\Aut_{\F_4}(\mathfrak l)\).
Their lower central dimensions agree after multiplication by two.
\end{theorem}

\begin{proof}
Lemma~\ref{lem:operator-pencil} provides the preliminary actions.
We first normalize a section and separate the scalar weights of the
cocycle (Steps 1--2). The decisive step is to combine binary degree two
with intrinsic generation to force square zero (Step 3).
We then verify sufficiency, construct the auxiliary Lie algebra and
compare lower central series (Steps 4--5). Finally, tracking changes of
coordinates upgrades the normal form to an equivalence on morphisms
(Step 6).

\emph{Step 1: a normalized section.}
Choose an \(\F_4\)-basis \(e_i\) of \(U\), with arbitrary lifts \(v_i\)
in \(V\). Define a binary-linear section by
\(e_i\mapsto v_i\), \(\omega e_i\mapsto f(v_i)\).
The images form a basis modulo \(W\).
Since \(f(v_i+f(v_i))=v_i\), the cocycle \(q\) in
\eqref{eq:preliminary-form} vanishes on each coordinate \(\F_4\)-line.

\emph{Step 2: scalar weights.}
Initially put \(P_i=L_{e_i}\), without yet asserting \(P_i^2=0\).
Lemma~\ref{lem:operator-pencil} gives
\[
 L_{\alpha e_i}=\alpha P_i+(\alpha+\alpha^2)P_i^2
       \quad(\alpha\in\F_4).
\]
Set \(N(x)=\sum_i x_i(P_i+P_i^2)\).
Then \(L_x=N+N^2\), \(N^3=0\), and
\(L_{\omega x}=\omega N+\omega^2N^2\).
Decompose \(q=q_0+q_1+q_2\) into the three scalar weights
\(q_j(\omega x)=\omega^jq_j(x)\). Projection to each weight is a linear
combination of \(q(x),q(\omega x),q(\omega^2x)\), so preserves binary
degree and the vanishing on coordinate lines.

The cocycle equation is
\[
 q(\omega x)=\omega^2q(x)+(\omega N+\omega^2N^2)q(x).
\]
Its weight components give
\[
 q_0=Nq_2+\omega N^2q_1,\qquad
 q_1=\omega Nq_0+\omega^2N^2q_2.
\]
Using \(N^3=0\) yields
\begin{equation}\label{eq:weight-solution}
 q_0=Nq_2,\qquad q_1=N^2q_2.
\end{equation}

\emph{Step 3: binary degree forces square zero.}
The reduced \(\F_4\)-polynomial of a weight-two binary-quadratic map has
the form \(\sum_i b_ix_i^2+\sum_{i<j}a_{ij}x_ix_j\).
Line normalization removes the first sum, so \(q_2=Q\).
Consequently
\[
 q_1=N^2Q=\sum_{k,i<j}x_k^2x_ix_jP_k^2a_{ij}.
\]
This is where the degree bound becomes decisive. For \(k=i\) or \(j\),
the exponent pattern is \((3,1)\); otherwise it is \((2,1,1)\).
Both patterns have binary degree three, whereas the weight projection
\(q_1\) has degree at most two. The cubed or squared coordinate
identifies \(k\), and the remaining coordinates identify \(\{i,j\}\).
Thus each displayed term has its own reduced monomial, forcing
\begin{equation}\label{eq:squares-kill-coefficients}
 P_k^2a_{ij}=0\qquad\text{for all }k,i,j.
\end{equation}
For completeness, the one-coordinate function basis
\(1,x,x^2,x^3\) has binary degrees \(0,1,1,2\).
Its leading classes form a basis in the associated graded degree
filtration. Tensoring over the coordinates proves that distinct reduced
monomials of the displayed degree three cannot cancel to degree two.
This justifies the comparison in the binary degree filtration.

We now extend this vanishing from the cocycle coefficients to the whole
intrinsic derived ideal. Let \(\m W=\sum_iP_iW\). By the preliminary identities,
every \(P_k^2\) annihilates \(\m W\).
Modulo \(\m W\), the polar values in
\eqref{eq:preliminary-form} are exactly the values of the polar of \(Q\).
Since \(W\) is the span of all polar values, this proves
\eqref{eq:intrinsic-generation}.
Together with \eqref{eq:squares-kill-coefficients}, it gives \(P_k^2=0\)
on all of \(W\). Now \(L_x=P(x)\) and \(q=Q+PQ\).
The degree-three terms of \(PQ\) with three distinct indices vanish
precisely when \eqref{eq:triangle} holds. The other terms have exponent
patterns \((2,1)\), of binary degree two. This proves necessity.

\emph{Step 4: sufficiency.}
Conditions \eqref{eq:square-zero} and \eqref{eq:triangle} make
\(P(x)^2=0\) and make \(PQ\) binary-quadratic.
Under \(x\mapsto\omega x\), the functions \(Q\) and \(PQ\) have weights
two and zero. Thus
\[
 (Q+PQ)(\omega x)=\omega^2Q+PQ
   =\omega^2(Q+PQ)+\omega P(Q+PQ).
\]
Together with
\((\omega\id+\omega P)(\omega\id+P)=\omega^2\id+P\),
this verifies \(f^2=f+\id\) by substitution.
Mixed polar values generate \(\m W\); modulo that space, the values
of the polar of \(Q\) generate the \(a_{ij}\)'s over \(\F_4\).
Equation~\eqref{eq:intrinsic-generation} therefore makes the polar
image exactly \(W\).

\emph{Step 5: the associated Lie algebra.}
Put \(\beta(x,y)=Q(x+y)+Q(x)+Q(y)\) and define
\begin{equation}\label{eq:auxiliary-bracket}
 [(x,z),(y,t)]_0=(0,P(x)t+P(y)z+\beta(x,y)).
\end{equation}
Jacobi follows from commuting operators and \eqref{eq:triangle}.
This is a metabelian \(\F_4\)-Lie algebra with derived ideal \(W\),
and it satisfies \eqref{eq:Lie-action-square}.
Conversely, splitting the abelianization of any algebra with these
properties gives exactly such operators and an alternating cocycle.
Choose its cross-term refinement \(Q\) to recover
\eqref{eq:normal-form}.

For both the original polar bracket and \eqref{eq:auxiliary-bracket},
the action on \(W\) is \(P(x)\). Hence, with \(\m\) the ideal generated
by these actions,
\begin{equation}\label{eq:lower-central-module}
 \gamma_2=W,\qquad\gamma_k=\m^{k-2}W\quad(k\ge3).
\end{equation}
This proves the lower central comparison.

\emph{Step 6: morphisms.}
A binary-linear conjugacy preserves the intrinsic space \(W\).
Its diagonal blocks \(A:U\to U'\) and \(C:W\to W'\) commute with
the intrinsic scalar actions, hence are \(\F_4\)-linear.
Write its off-diagonal block as \(t_s+t_a\), where
\(t_s(\omega x)=\omega t_s(x)\) and
\(t_a(\omega x)=\omega^2t_a(x)\).
Direct substitution in \(Sf=f'S\) gives the two conditions
\begin{align}
 CP(x)&=P'(Ax)C,\label{eq:morphism-action}\\
 Q'(Ax)&=CQ(x)+t_a(x)+P'(Ax)t_s(x).
 \label{eq:morphism-cocycle}
\end{align}
Indeed the unseparated cocycle equation reduces to the second one by
multiplication with \(\id+P'(Ax)\), whose inverse is itself.
Polarization over \(\F_4\) eliminates \(t_a\) and gives
\begin{equation}\label{eq:morphism-Lie}
 \beta'(Ax,Ay)=C\beta(x,y)
              +P'(Ax)t_s(y)+P'(Ay)t_s(x).
\end{equation}
Equations~\eqref{eq:morphism-action} and \eqref{eq:morphism-Lie}
say exactly that \((x,z)\mapsto(Ax,Cz+t_s(x))\) is an
\(\F_4\)-Lie isomorphism.

Conversely, for any such Lie isomorphism the function
\(Q'(Ax)+CQ(x)+P'(Ax)t_s(x)\) has zero polar and is an ordinary
homogeneous quadratic. It is therefore a unique sum of squares
\(\sum b_ix_i^2\). Taking this as \(t_a\) gives a unique conjugacy.
Scalar parts respect composition because \(A\) and \(C\) are
\(\F_4\)-linear. Thus the correspondence is an equivalence on both
objects and morphisms, proving also the centralizer assertion.
\end{proof}

\begin{corollary}\label{cor:metabelian-nilpotent}
Every metabelian polar algebra is nilpotent and satisfies
\((\ad v)^3=0\) for every \(v\). If \(w=\dim_{\F_4}W\), its class is
at most \(w+1\).
\end{corollary}
\begin{proof}
The commuting square-zero operators generate a nilpotent associative ideal.
They have a common nonzero kernel whenever \(W\ne0\); repeating on the
quotient gives a strictly triangular flag, so all products of \(w\)
operators vanish. Apply \eqref{eq:lower-central-module}.
Finally \(\ad v\) maps \(V\) into \(W\), and its square on \(W\) is zero.
\end{proof}

\begin{theorem}\label{thm:solvable-nilpotent}
A polar Lie algebra is solvable if and only if it is nilpotent.
Its solvable radical is invariant under \(f\) and is nilpotent.
\end{theorem}
\begin{proof}
Induct on derived length. The abelian case is immediate.
Put \(W=L'\), \(W'=L''\). The restricted solution on \(W\) has
nilpotent polar algebra by induction. The quotient \(L/W'\) is
metabelian, hence nilpotent by Corollary~\ref{cor:metabelian-nilpotent}.
For fixed \(x\in L\), the derivation \(D=\ad x|_W\) is therefore
nilpotent on \(W/W'\).

We use the following elementary derivation fact. If \(N\) is nilpotent
and \(D\) is nilpotent on \(N/N'\), then it is nilpotent on \(N\).
Each factor \(\gamma_i(N)/\gamma_{i+1}(N)\) is a quotient of
\((N/N')^{\otimes i}\). The induced derivation is the sum of the \(i\)
commuting operators acting in individual tensor positions.
If \(D^s=0\) on \(N/N'\), the sum has power
\(i(s-1)+1\) equal to zero, by the pigeonhole principle on its expanded
products. The finite lower central filtration proves the assertion on
\(N\), in every characteristic.

Apply this fact to \(N=W\). Since \(\ad x\) maps \(L\) into \(W\),
it is nilpotent on \(L\). Engel's theorem now gives nilpotency of \(L\).
The converse is immediate. If \(R\) is the solvable radical,
\(L/R\) has zero center. Lemma~\ref{lem:ideal-descent} makes \(R\)
invariant, and the result just proved applies to the restricted solution.
\end{proof}

\section{The universal module and sharp metabelian bounds}
\label{sec:universal}

\subsection{A multigraded commutator module}

For \(d\ge2\), put
\[
 R_d=\F_4[t_1,\ldots,t_d]/(t_1^2,\ldots,t_d^2),
 \qquad\m=(t_1,\ldots,t_d).
\]
Let \(W_d\) be the \(R_d\)-module generated by \(v_{ij}\), \(i<j\),
subject to
\begin{equation}\label{eq:module-triangle}
 t_iv_{jk}+t_jv_{ik}+t_kv_{ij}=0\qquad(i<j<k).
\end{equation}
We set \(v_{ji}=v_{ij}\) and \(v_{ii}=0\).
This is the classical free metabelian commutator-module presentation,
base-changed to \(R_d\); compare~\cite[Theorem~6.2]{PapadimaSuciu2004}
and~\cite[Sections~2--3]{KovacsStohr2013}.

Every derived module in Theorem~\ref{thm:correspondence} is a quotient of
\(W_d\): send \(t_i\) to \(P_i\) and \(v_{ij}\) to \(a_{ij}\).
The intrinsic-generation condition makes this homomorphism surjective.
Conversely, every quotient of \(W_d\), with these images, gives the
normal form \eqref{eq:normal-form}.

Grade \(R_d\) by coefficient degree, and put the generators \(v_{ij}\)
in coefficient degree zero. We also use the finer grading
\(\deg t_i=e_i\), \(\deg v_{ij}=e_i+e_j\).

\begin{proposition}\label{prop:Hilbert}
The coefficient-degree Hilbert polynomial of \(W_d\) is
\begin{equation}\label{eq:Hilbert}
 H_d(T)=\sum_{k=2}^d\binom dk
       \bigl((k-1)T^{k-2}+kT^{k-1}\bigr)+\binom d2T^2.
\end{equation}
Consequently
\[
 \dim_{\F_4}W_d=(d-1)2^d+1-d+\binom d2.
\]
The universal solution on \(\F_4^d\oplus W_d\) has class
\(\max(4,d+1)\).
\end{proposition}
\begin{proof}
A raw multihomogeneous monomial is \(t_Sv_{ij}\), where \(S\) is
square-free. Its multidegree has entries zero, one, or two, and at most
two entries are doubled. Fix its support of size \(k\ge2\).

If no index is doubled, the raw coordinates are the edges of the complete
graph on the support. Relations~\eqref{eq:module-triangle} are its triangle
sums. The quotient edge space has dimension \(k-1\), in coefficient
degree \(k-2\).

If exactly one index \(j\) is doubled, only edges incident with \(j\)
occur. Triangle relations through \(j\) identify every pair of these
edges, leaving dimension one. There are \(k\) choices of \(j\), in
coefficient degree \(k-1\).

If exactly two indices \(i,j\) are doubled, there is one raw coordinate.
A third index in the support kills it by a triangle relation, since the
other two terms contain a squared variable. The coordinate survives
exactly when the support is \(\{i,j\}\), giving the final term of
\eqref{eq:Hilbert}. These cases exhaust the multigrading.

The maximum coefficient degree is two for \(d=2,3\), and \(d-1\) for
\(d\ge4\). Equation~\eqref{eq:lower-central-module} gives the class.
\end{proof}

In particular, binary degree two permits unbounded polar nilpotency class.
For the exact minimum dimension we must retain a terminal commutator in
as small a quotient of \(W_d\) as possible.

\subsection{Terminal tensors and contraction modules}

\begin{lemma}\label{lem:terminal-repeat}
Let \(\mathfrak l\) be a metabelian \(\F_4\)-Lie algebra satisfying
\eqref{eq:Lie-action-square}, of class \(c\ge3\).
Some nonzero terminal commutator has a repeated input between its first
bracket pair and an operator position. If \(c\ge5\), there is a
subalgebra of the same class generated by exactly \(c-1\) elements
modulo its derived ideal.
\end{lemma}
\begin{proof}
The last lower central ideal is spanned by
\[
 T(x_1,x_2;x_3,\ldots,x_c)
     =P(x_3)\cdots P(x_c)\beta(x_1,x_2).
\]
Indeed \(W=\Span\beta+\m W\), and multiplication by \(\m^{c-2}\)
kills the second summand. The operator positions are symmetric, and a
repeated input within them gives zero.

Suppose every value \(T(x,y;x,z_4,\ldots,z_c)\) were zero.
Polarizing the repeated input gives
\[
 T(x,y;z,\ldots)=T(z,y;x,\ldots).
\]
Together with symmetry of the first pair and the operator positions,
this makes \(T\) fully symmetric. Jacobi on the first three inputs
then gives \(3T=T=0\), a contradiction.

Choose a nonzero repeated-input value. Its inputs span a subspace
\(U_0\) of the abelianization of dimension at most \(c-1\).
Their chosen lifts generate a subalgebra \(H\) of class \(c\), with
abelianization exactly \(U_0\): all words of length at least two lie
in its derived ideal. Its module is a quotient of \(W_{d_0}\), where
\(d_0=\dim U_0\).
By Proposition~\ref{prop:Hilbert}, class \(c\ge5\) requires
\(d_0\ge c-1\). Thus \(d_0=c-1\).
\end{proof}

For \(d\ge4\), the top coefficient degree of \(W_d\) is \(d-1\).
It has basis \(T_1,\ldots,T_d\), where
\begin{equation}\label{eq:terminal-basis}
 T_j=\Bigl(\prod_{h\ne i}t_h\Bigr)v_{ij}\quad(i\ne j).
\end{equation}
Triangle relations identify the choices of \(i\).

\begin{lemma}\label{lem:top-contraction}
The top component is the \(\GL(d,4)\)-module \(\det(U)\otimes U\).
For every nonzero functional \(\lambda\) on that component, extended
by zero to the other degrees,
\[
 \dim_{\F_4}R_d\lambda=2^d-2.
\]
If a quotient \(M\) of \(W_d\) retains a nonzero top component, then
\(\dim_{\F_4}M\ge2^d-2\), whether or not the quotient is graded.
\end{lemma}
\begin{proof}
Under a diagonal basis change \(A\), \(T_j\) scales by
\(\det(A)A_{jj}\). A transvection \(e_j\mapsto e_j+\alpha e_k\)
sends \(T_j\) to \(T_j+\alpha T_k\), and fixes the other top basis
vectors. To check the formula choose \(i\notin\{j,k\}\) in
\eqref{eq:terminal-basis}; any extra product has a squared variable.
For the remaining vectors an index outside the affected entries can
be chosen since \(d\ge4\).
Together with permutations these transformations prove the module assertion.
They act transitively on nonzero dual vectors, including scalar multiples.

Normalize \(\lambda(T_j)=1\), \(\lambda(T_k)=0\) for \(k\ne j\).
We claim
\begin{equation}\label{eq:annihilator}
 \Ann_{R_d}(\lambda)=\Bigl(\prod_{i\ne j}t_i\Bigr).
\end{equation}
The annihilator is multigraded: contractions by distinct monomials have
distinct resulting dual multidegrees. For a square-free monomial \(t_S\)
not containing \(j\), the remaining multidegree is of the one-doubled-index
type. It survives unless \(S\) contains every other index.
If \(j\in S\), the remaining multidegree has no doubled index.
It survives whenever an index is omitted, as witnessed by an edge
joining that index to \(j\); it vanishes when all indices are used.
These are exactly the two monomials in the ideal on the right of
\eqref{eq:annihilator}. This proves the claim and the dimension formula.

For an arbitrary quotient \(M\), pull back a functional nonzero on the
top image to \(\phi\in W_d^*\). Its highest homogeneous component
\(\phi_{d-1}\) is nonzero. Independent homogeneous contractions of
\(\phi_{d-1}\) lift to independent contractions of \(\phi\), by
comparison of their leading dual degrees. Hence
\[
 \dim M\ge\dim R_d\phi\ge\dim R_d\phi_{d-1}=2^d-2.
\]
This is the usual filtered contraction argument; it does not assume a
graded kernel for \(W_d\to M\).
\end{proof}

\subsection{Exact minima}

\begin{theorem}\label{thm:unrestricted-minimum}
Among metabelian quadratic solutions of polar class \(c\ge2\), the
minimum binary dimension is
\[
 \mu(2)=6,\quad\mu(3)=8,\quad\mu(4)=12,\qquad
 \mu(c)=2^c+2c-6\quad(c\ge5).
\]
\end{theorem}
\begin{proof}
Write \(d=\dim_{\F_4}(V/W)\), \(w=\dim_{\F_4}W\).
Nonabelian objects have \(d\ge2\) by Lemma~\ref{lem:two-generators}.
A nonzero product \(P_1\cdots P_kv\), for commuting square-zero
operators, yields \(2^k\) independent vectors
\[
 v_S=\Bigl(\prod_{i\in S}P_i\Bigr)v\qquad(S\subseteq\{1,\ldots,k\}).
\]
In a relation, choose a term with nonzero coefficient of smallest support
and multiply by its complementary operators. Every other surviving
term repeats an operator and vanishes, while the chosen term does not.
Taking \(k=c-2\) gives \(w\ge2^{c-2}\).
For \(c=2,3,4\) this gives the respective lower bounds \(6,8,12\).

They are attained by cyclic quotients of \(W_2=R_2v_{12}\): use
\(R_2/\m\), \(R_2/(t_2)\), and \(R_2\), respectively, as the
derived module in \eqref{eq:normal-form}.
The first two representatives are
\begin{align*}
 &(a,b,z)\longmapsto(\omega a,\omega b,\omega z+ab),\\
 &(a,b,y,z)\longmapsto
    (\omega a,\omega b,\omega y+az+a^2b,\omega z+ab).
\end{align*}
The class-four representative appears in
\eqref{eq:seven-coordinates} below.

Let \(c\ge5\). Proposition~\ref{prop:Hilbert} gives \(d\ge c-1\).
Lemma~\ref{lem:terminal-repeat} supplies a same-class subalgebra with
\(d_0=c-1\). Its derived module retains the top of \(W_{c-1}\), so
Lemma~\ref{lem:top-contraction} gives
\(w\ge2^{c-1}-2\). Therefore
\[
 \dim_{\F_2}V=2(d+w)\ge2^c+2c-6.
\]

To attain the bound, set \(d=c-1\ge4\), choose a top functional
\(\lambda\) as in Lemma~\ref{lem:top-contraction}, and put
\begin{equation}\label{eq:terminal-quotient}
 M_d=W_d/\{w:\lambda(rw)=0\text{ for every }r\in R_d\}.
\end{equation}
The denominator is an \(R_d\)-submodule. The quotient has dimension
\(2^d-2\) and retains the terminal value, so its solution has class
\(d+1=c\) and binary dimension \(2(d+2^d-2)\).
\end{proof}

The change of formula at class five comes from the top degree of the
universal module. For two and three quotient generators,
\[
 H_2(T)=1+2T+T^2,\qquad H_3(T)=3+8T+6T^2.
\]
In these cases the degree-two double-endpoint terms contribute at the
top, allowing class four even with two quotient generators. For
\(d\ge4\), they lie below the top degree \(d-1\); the terminal component
\(\det(U)\otimes U\) then controls class \(d+1\) uniformly. This is why
the contraction argument yields one formula from class five onward,
while the cyclic quotients of \(W_2\) give the three smaller values.

The extremal quotient \(M_d\) in \eqref{eq:terminal-quotient} has
coefficient Hilbert polynomial
\[
 (d-1)+\sum_{k=1}^{d-1}\binom d{d-1-k}T^k.
\]
For example, the class-five extremizer has derived Hilbert vector
\((3,6,4,1)\) and binary lower central dimensions
\((36,28,22,10,2,0)\).
The contraction construction supplies an explicit finite module
description at every class, without an orbit search.

\section{When the original polar bracket is \texorpdfstring{\(\F_4\)}{F4}-bilinear}
\label{sec:compatible}

Return to the normal form \eqref{eq:normal-form} and keep the intrinsic
scalar actions on \(U=V/W\) and \(W\). Its auxiliary bracket
\eqref{eq:auxiliary-bracket} is always \(\F_4\)-bilinear.
The following definition concerns the original polar bracket of \(f\).

\begin{definition}\label{def:compatible}
A metabelian solution is \emph{compatible} if \(V\) admits an
\(\F_4\)-vector space structure that restricts to the intrinsic structure
on \(W\), induces the intrinsic structure on \(U\), and makes the
original polar bracket \(\F_4\)-bilinear.
\end{definition}

\begin{theorem}\label{thm:compatibility}
For a normal form with data \(P_i,a_{ij}\), the following are equivalent:
\begin{enumerate}
\item the solution is compatible;
\item there are \(b_1,\ldots,b_d\in W\) satisfying
\begin{equation}\label{eq:compatibility-system}
 P_ib_i=0,\qquad P_ib_j=P_ja_{ij}\quad(i\ne j);
\end{equation}
\item the auxiliary \(\F_4\)-Lie algebra has a restricted square \(r\)
with \(r(W)=0\) and \(r(\mathfrak l)\subseteq W\).
\end{enumerate}
If a compatible scalar lift exists, the number of such scalar actions
on the fixed binary space, with the prescribed actions on \(U,W\), is
\(4^{dk}\), where \(k=\dim_{\F_4}\bigcap_i\ker P_i\).
\end{theorem}
\begin{proof}
Every binary-linear lift of multiplication by \(\omega\) has the form
\[
 T(x,z)=(\omega x,\omega z+t(x)).
\]
The relation \(T^2+T=\id\) is exactly
\(t(\omega x)=\omega^2t(x)\). Thus
\(t(x)=b(x)=\sum_i x_i^2b_i\).
The involutory binary shear \(S(x,z)=(x,z+b(x))\) conjugates \(T\)
to standard scalar multiplication. It changes the cocycle of \(f\) to
\[
 Q+b+P(Q+b).
\]
The first two terms have \(\F_4\)-bilinear polars. The last term has
scalar weight zero. Its polar could be \(\F_4\)-bilinear only if that
polar were zero, since an \(\F_4\)-bilinear polar has weight two.
An additive zero-constant function of weight zero is zero: additivity
and invariance under \(\omega\) give its value at \(\omega^2x=x+\omega x\)
equal to zero, and then invariance gives its value at \(x\) equal to zero.
Compatibility is therefore equivalent to
\begin{equation}\label{eq:compatibility-polynomial}
 P(x)(Q(x)+b(x))=0.
\end{equation}
After the triangle cancellations, comparison of \(x_i^3\) and
\(x_ix_j^2\) gives precisely \eqref{eq:compatibility-system}.

We recall explicitly the restricted-square axioms in characteristic two:
\(r(\alpha u)=\alpha^2r(u)\),
\(r(u+v)=r(u)+r(v)+[u,v]_0\), and
\(\ad r(u)=(\ad u)^2\).
Any square with the range and vanishing conditions in (iii) must have
the form
\begin{equation}\label{eq:restricted-square}
 r(x,z)=(0,P(x)z+q(x)),\qquad q(x)=Q(x)+b(x).
\end{equation}
Its first two axioms are automatic. For \(u=(x,z),v=(y,t)\), the two
adjoint expressions have final coordinates
\[
 P(x)P(y)z+P(x)\beta(x,y),\qquad
 P(x)P(y)z+P(y)q(x),
\]
respectively. Condition~\eqref{eq:compatibility-polynomial} implies
\(P(x)\beta(x,y)=P(y)q(x)\): take the coefficient of \(s\) in
\(P(x+sy)q(x+sy)=0\). Its degree in \(s\) is at most three, below
\(|\F_4|\), so the comparison is valid for functions over \(\F_4\).
This proves the third restricted axiom. Conversely that axiom, applied
to \(u\), gives \([u,r(u)]_0=0\), hence
\eqref{eq:compatibility-polynomial}.

Finally two solutions of \eqref{eq:compatibility-system} differ by
independent choices of \(d\) vectors in \(\bigcap_i\ker P_i\).
The displayed count follows. It counts scalar actions, not their
conjugacy classes.
\end{proof}

\begin{example}\label{ex:incompatible-eight}
The eight-dimensional class-three model has
\(P_1z=y\), \(P_1y=0\), \(P_2=0\), and \(a_{12}=z\).
Condition~\eqref{eq:compatibility-system} would require
\(0=P_2b_1=P_1a_{12}=y\). Hence it has no compatible scalar action.
This gives a concrete distinction between the auxiliary algebra and
the original polar algebra.
\end{example}

\begin{remark}
The restricted square in Theorem~\ref{thm:compatibility} is a square on
the auxiliary algebra. The theorem imposes both its range and its
restriction to the derived ideal. General correspondences between
quadratic maps and restricted Lie structures, such as
\cite[Section~9]{Pakianathan2007}, carry their own additional identities.
Here the exact additional condition is the solvable linear system
\eqref{eq:compatibility-system}.
\end{remark}

\section{Sharp compatible bounds and extremizers}
\label{sec:compatible-bounds}

\subsection{The abelianization bound}

\begin{lemma}\label{lem:compatible-d}
A nonabelian compatible solution of class \(c\), with
\(d=\dim_{\F_4}(V/W)\), satisfies \(c\le d\).
\end{lemma}
\begin{proof}
Use a restricted square from Theorem~\ref{thm:compatibility}.
The identity
\(P(x)\beta(x,y)=P(y)r(x)\), with a lift of \(x\) understood,
implies
\begin{equation}\label{eq:double-endpoint}
 P(x)P(y)\beta(x,y)=0.
\end{equation}
In the universal module this kills the submodule \(K\) spanned by
\(t_it_jv_{ij}\), \(i<j\).
It has dimension \(\binom d2\) and lies in coefficient degree two.
Every multiplication by an endpoint kills it by a square, and
multiplication by a third index kills it by a triangle relation.

This relation space is invariant under \(\GL(U)\).
Indeed the universal expression
\(P(x)P(y)\beta(x,y)\) is additive in each input and scales
quadratically in each. For example its additive cross term in \(x\)
is
\[
 P(y)\bigl(P(x)\beta(z,y)+P(z)\beta(x,y)\bigr)
       =P(y)^2\beta(x,z)=0
\]
by Jacobi. Its basis values therefore span precisely \(K\).

Put \(\overline W_d=W_d/K\). The multigrading in
Proposition~\ref{prop:Hilbert} shows that its highest coefficient degree
is \(d-1\), for every \(d\ge2\), and that its top has the basis
\(T_j\) of \eqref{eq:terminal-basis}.
This top is \(\det(U)\otimes U\) also for \(d=2,3\): the same
diagonal and transvection calculations apply, with any additional
double-endpoint term now in \(K\). For \(d=2\) this is immediate
from \(T_1=t_1v_{12},T_2=t_2v_{12}\).

Every compatible derived module \(M\) is a quotient of \(\overline W_d\),
so \(c\le d+1\). Suppose \(c=d+1\). Choose a functional on \(M\)
nonzero on its terminal component. A change of quotient basis normalizes
its top pullback to \(\lambda(T_j)=1\), \(\lambda(T_k)=0\) for
\(k\ne j\).
Using \(P_j a_{ij}=P_i b_j\), we have, for \(i\ne j\),
\begin{equation}\label{eq:terminal-square-value}
 \im(T_j)=\Bigl(\prod_{h\ne j}P_h\Bigr)b_j.
\end{equation}
Write \(b_j=\sum c_{kl}a_{kl}\pmod{\m M}\), using intrinsic
generation. The product on the right kills \(\m M\).
If \(j\) is an endpoint of \(v_{kl}\), then
\((\prod_{h\ne j}t_h)v_{kl}\) is a top vector whose doubled
endpoint is different from \(j\). If \(j\) is not an endpoint,
the vector has both endpoints doubled and is in \(K\).
Thus the normalized top functional is zero on the right side of
\eqref{eq:terminal-square-value}, but one on the left.
This contradiction proves \(c\le d\).
\end{proof}

\subsection{Counting preimages of an operator cube}

\begin{lemma}\label{lem:compatible-w}
A compatible solution of class \(c\ge2\) satisfies
\[
 \dim_{\F_4}W\ge2^{c-1}-1.
\]
\end{lemma}
\begin{proof}
The case \(c=2\) says \(W\ne0\). Let \(c\ge3\).
By Lemma~\ref{lem:terminal-repeat}, choose a nonzero terminal value
\[
 P(x)P(z_1)\cdots P(z_{c-3})\beta(x,y).
\]
Restrictedness gives \(P(x)\beta(x,y)=P(y)b\), where
\(b=r(x)\in W\). Put \(m=c-2\),
\(u_1=y,u_2=z_1,\ldots,u_m=z_{c-3}\).
Then \(P(u_1)\cdots P(u_m)b\ne0\).
The \(2^m\) vectors
\[
 b_S=\Bigl(\prod_{i\in S}P(u_i)\Bigr)b
 \qquad(S\subseteq\{1,\ldots,m\})
\]
are independent by the complementary-operator argument used in
Theorem~\ref{thm:unrestricted-minimum}. Let \(C\) be their span.
Since \([x,r(x)]_0=0\), \(P(x)\) kills \(b\), hence all of \(C\).

For each nonempty \(S\), choose \(i\in S\) and set
\[
 v_S=\Bigl(\prod_{j\in S\setminus\{i\}}P(u_j)\Bigr)\beta(x,u_i).
\]
These vectors lie in the intrinsic derived ideal, and restrictedness gives
\(P(x)v_S=b_S\). Thus
\[
 \dim\ker(P(x)|_W)\ge2^m,\qquad
 \rank(P(x)|_W)\ge2^m-1.
\]
Rank-nullity yields \(\dim W\ge2^{m+1}-1=2^{c-1}-1\).
\end{proof}

\begin{theorem}\label{thm:compatible-minimum}
For every \(c\ge2\), the minimum binary dimension among compatible
solutions of polar class \(c\) is
\[
 \nu(c)=2^c+2c-2.
\]
\end{theorem}
\begin{proof}
Lemmas~\ref{lem:compatible-d} and \ref{lem:compatible-w} give the
lower bound by addition.
For attainment put \(s=c-1\),
\(R_s=\F_4[t_1,\ldots,t_s]/(t_i^2)\), and use
\begin{equation}\label{eq:semidirect-benchmark}
 f(x,z)=\left(\omega x,\omega z+\Bigl(\sum_i x_it_i\Bigr)z\right)
       \quad\text{on }\F_4^s\oplus R_s.
\end{equation}
The square of \(\sum x_it_i\) is zero, so direct substitution gives
\eqref{eq:functional}. The polar bracket is \(\F_4\)-bilinear and
has intrinsic derived ideal \(\m R_s\).
The scalar actions induced by \(f\) on that ideal and on the
abelianization are the displayed ones, so the map is compatible.
Its lower central ideals are
\(\gamma_k=\m^{k-1}R_s\) for \(k\ge2\), giving class \(s+1=c\).
The derived and quotient dimensions are \(2^s-1\) and \(s+1\).
The binary dimension is therefore \(2(s+2^s)=2^c+2c-2\).
\end{proof}

For \(c\ge5\), the extra four binary dimensions have an exact structural
meaning. The sharp bounds in the two subclasses are
\[
\begin{array}{c|cc}
 & d & w\\ \hline
\text{metabelian} & c-1 & 2^{c-1}-2\\
\text{compatible metabelian} & c & 2^{c-1}-1
\end{array}
\]
where each entry is the corresponding lower bound, attained
simultaneously. Compatibility therefore requires one additional
\(\F_4\)-dimension in the quotient and one in the derived ideal,
giving four additional binary dimensions. The restricted square
explains the strengthened bounds: it constrains terminal commutators
and supplies independent preimages for the nonempty faces of the
operator cube. The elementary semidirect construction attains this cost.

\subsection{Several extremizers in every class}

\begin{theorem}\label{thm:extremizers}
For every \(c\ge4\), at least \(c-2\) binary-linear conjugacy classes
attain \(\nu(c)\).
\end{theorem}
\begin{proof}
Set \(s=c-1\ge3\) and
\[
 R=\F_4[t_0,\ldots,t_{s-1}]/(t_i^2),\qquad W=\m R.
\]
Let \(U\) have basis \(x,u_1,\ldots,u_{s-1},v\).
Choose either \(h=0\), or \(h=t_1\cdots t_r\) with
\(1\le r\le s-2\), and put \(p=t_0h\).
Define the action on \(W\) by multiplication:
\[
 P_x=t_0,\qquad P_{u_i}=t_i,\qquad P_v=p.
\]
Define the alternating cocycle by
\[
 \beta(x,u_i)=t_i,\quad\beta(x,v)=t_0,\quad
 \beta(v,u_i)=(1+h)t_i,\quad\beta(u_i,u_j)=0.
\]
The only Jacobi triple types to check are \((x,u_i,u_j)\),
\((v,u_i,u_j)\), and \((x,v,u_i)\). Their sums are, respectively,
\[
 t_it_j+t_jt_i,\quad
 t_i(1+h)t_j+t_j(1+h)t_i,\quad
 t_0(1+h)t_i+pt_i+t_it_0,
\]
all zero. The cocycle values generate \(W\) under the action.

Take square values \(b_x=t_0,b_{u_i}=0,b_v=p\). They satisfy
\eqref{eq:compatibility-system}; the nontrivial relation involving \(p\)
uses \(p(1+h)=p\), since \(h^2=0\).
In scalar coordinates \(a=(X,Y_1,\ldots,Y_{s-1},V)\), the resulting
restricted square has the form
\begin{align}
 r(a,w)&=(0,P(a)w+q(a)),\label{eq:deformation-square}\\
 q(a)&=t_0(X^2+XV+hV^2)
       +\sum_i t_iY_i\bigl(X+(1+h)V\bigr).
       \label{eq:deformation-q}
\end{align}
All ring elements in these formulas are fixed coefficients. Consequently
\(f=\omega\id+r\) has binary degree at most two.
The square vanishes on \(W\), takes values in \(W\), and
\([a,r(a)]_0=0\). Expanding \(r(\omega a+r(a))\) proves
\(f^2=f+\id\). Its original polar is the displayed \(\F_4\)-bracket,
so it is compatible with the intrinsic scalar actions.

The lower central ideals are \(\gamma_k=\m^{k-1}R\) for \(k\ge2\),
since the action includes every \(t_i\). Each map has class \(c\)
and dimension \(\nu(c)\).

To distinguish them, use only intrinsic binary invariants. Let
\[
 \mathcal A=\im\bigl(L\longrightarrow\End_{\F_2}(W),\
                   a\longmapsto\ad a|_W\bigr),\qquad
 \mathcal A_k=\{D\in\mathcal A:D(W)\subseteq\gamma_{k+2}(L)\}.
\]
For \(h=0\), \(\dim_{\F_4}\mathcal A=s\) and
\(\mathcal A_k=0\) for \(k\ge2\).
For \(h=t_1\cdots t_r\), multiplication by \(p\) is independent of
the linear multiplications and is nonzero on \(W\), as \(r\le s-2\).
Thus \(\dim_{\F_4}\mathcal A=s+1\), and
\[
 \dim_{\F_4}\mathcal A_k=
 \begin{cases}1,&2\le k\le r+1,\\0,&k\ge r+2.\end{cases}
\]
Indeed any nonzero linear multiplication has a nonzero degree-two
output on \(W\), whereas \(pW\) has minimum nonzero degree \(r+2\).
The binary dimensions are twice these dimensions and are preserved by
every binary-linear conjugacy, independently of a choice of scalar lift.
They separate \(h=0\) and the \(s-2\) possible values of \(r\).
\end{proof}

The theorem gives an explicit collection of extremizers. It leaves open
the classification of all equality cases and their automorphism groups.

\section{Two closed classification ranges}
\label{sec:classification}

The polar framework supports two complementary classification questions:
fix the ambient dimension, or fix the abelianization while allowing the
ambient dimension to grow. The first application combines simple-algebra
obstructions with the metabelian correspondence to close dimension eight.
The second uses central descent and exact extension spaces to close the
entire solvable abelianization-four stratum. Its last two members are
nonmetabelian, so the extension method also reaches beyond the
correspondence of Section~\ref{sec:metabelian}.

\subsection{Obstructions to simple polar algebras}

The simple reduction in Corollary~\ref{cor:simple-reduction} can be
combined with two explicit obstructions.

\begin{proposition}\label{prop:matrix-obstruction}
No polar algebra has a Lie subquotient isomorphic to the underlying
binary Lie algebra of \(\mathfrak{sl}_3(K)\), for a finite field \(K\)
of characteristic two. In particular it cannot be
\(\mathfrak{sl}_m(K)\) or \(\mathfrak{psl}_m(K)\), \(m\ge3\).
\end{proposition}
\begin{proof}
In \(\mathfrak{sl}_3(K)\), take
\(a=E_{12},b=E_{23},c=E_{31},d=E_{11}+E_{22}\).
The three terms in \eqref{eq:four-variable} are
\(E_{11}+E_{33},E_{22}+E_{33},0\), whose sum is nonzero.
A polynomial identity passes to subalgebras and quotients.
The upper \(3\times3\) block gives the required subalgebra in
\(\mathfrak{sl}_m(K)\). For \(m\ge4\) it also embeds in the quotient
by scalar matrices, since the block intersects that center trivially.
For \(m=3\) the center is zero.
\end{proof}

\begin{lemma}[Two centralizers]\label{lem:two-centralizers}
Let \(L\) be a Lie algebra over a characteristic-two field \(K\).
If \(x,y\in L\) satisfy
\[
 C_L(x)=Kx,\qquad C_L(x+y)=K(x+y),\qquad
 [y,[x,y]]\notin K[x,y],
\]
then there is no zero-constant binary-quadratic map with polar bracket
\([\ ,\ ]\) satisfying \([v,f(v)]=0\) for all \(v\).
\end{lemma}
\begin{proof}
The commuting condition forces
\(f(x)=\alpha x\), \(f(x+y)=\beta(x+y)\), for some \(\alpha,\beta\in K\).
Polarization then gives
\(f(y)=(\alpha+\beta)x+\beta y+[x,y]\).
Commuting this value with \(y\) contradicts the third hypothesis.
No \(K\)-linearity of \(f\) is used.
\end{proof}

\begin{proposition}\label{prop:Witt-obstruction}
For every finite characteristic-two field \(K\) and \(r\ge2\), the
derived divided-power Witt algebra \(W(1;r)'\) over \(K\), viewed over
\(\F_2\), is not a polar algebra of a quadratic solution.
\end{proposition}
\begin{proof}
Use the classical divided-power algebra \(A=O(1;r)\), with basis
\(x^{(a)}\), \(0\le a<2^r\), and multiplication
\[
 x^{(a)}x^{(b)}=\binom{a+b}{a}x^{(a+b)}
\]
when \(a+b<2^r\), and zero otherwise. Let
\(D(x^{(a)})=x^{(a-1)}\), with \(D(1)=0\).
Then \(\ker D=K\), and
\(W(1;r)'=\Span_K\{x^{(a)}D:0\le a\le2^r-2\}\);
see~\cite[Section~4.5]{LopatkinZusmanovich2019}.

If \(u\) is a unit of \(A\), a commuting vector \(vD\) satisfies
\(u^2D(v/u)=0\), so \(C_{AD}(uD)=KuD\).
Choose \(x=D,y=x^{(2)}D\). Both \(1\) and \(1+x^{(2)}\) are units,
so the first two hypotheses of Lemma~\ref{lem:two-centralizers} hold
also inside \(W(1;r)'\). Finally
\[
 [x,y]=x^{(1)}D,\qquad [y,[x,y]]=x^{(2)}D,
\]
because \((x^{(1)})^2=0\). These are \(K\)-independent.
\end{proof}

\subsection{All eight-dimensional solutions}

\begin{theorem}\label{thm:dimension-eight}
In binary dimension eight there are exactly three linear-conjugacy
classes of quadratic solutions. Representatives on \(\F_4^4\) are
\begin{align*}
 L(a,b,c,d)&=(\omega a,\omega b,\omega c,\omega d),\\
 C(a,b,c,d)&=(\omega a,\omega b,\omega c,\omega d+ab),\\
 N(a,b,y,z)&=(\omega a,\omega b,\omega y+az+a^2b,\omega z+ab).
\end{align*}
Their polar-image dimensions are \(0,2,4\), their center dimensions
are \(8,4,2\), and their classes are \(1,2,3\), respectively.
\end{theorem}
\begin{proof}
We first exclude nonsolvable polar algebras in every binary dimension
at most eight. Corollary~\ref{cor:simple-reduction} would give an
even-dimensional nonabelian simple polar algebra of dimension at most
eight. The classical lists of Strade~\cite{Strade2006} and
Vaughan-Lee~\cite[Theorem~1]{VaughanLee2006} leave the six-dimensional
algebra obtained from \(W(1;2)'\) over \(\F_4\), and the two
eight-dimensional simple types over \(\F_2\).
The first is excluded by Proposition~\ref{prop:Witt-obstruction}.

For the other two, use the published bases \(b_1,\ldots,b_8\) in
the two tables of~\cite[Section~3.6]{LieAlgDB}.
The three terms of \eqref{eq:four-variable}, at
\((b_1,b_2,b_3,b_4)\), are as follows:
\begin{center}
\begin{tabular}{@{}lccc@{}}
\toprule
Type&\([[b_1,b_2],[b_3,b_4]]\)&
\([[b_1,b_3],[b_2,b_4]]\)&\([[b_1,b_4],[b_2,b_3]]\)\\
\midrule
First&\(0\)&\(0\)&\(b_8\)\\
Second&\(b_2+b_3+b_4\)&\(b_2+b_3\)&\(b_2+b_4\)\\
\bottomrule
\end{tabular}
\end{center}
Both sums are nonzero. The list is an external classification input;
these short evaluations supply the equation-specific exclusions.

Thus all polar algebras in the stated dimensions are nilpotent by
Theorem~\ref{thm:solvable-nilpotent}. A nonabelian nilpotent polar
algebra has dimension at least six, by
Lemma~\ref{lem:two-generators} and parity of its nonzero derived ideal.
In dimension eight its derived ideal has dimension at most four and
is therefore abelian. The metabelian correspondence applies.

If \(W=0\), there is the single linear class.
If \(\dim_{\F_4}W=1\), all square-zero actions vanish; the cocycle
is an alternating scalar form on \(\F_4^3\). Its nonzero rank is two,
giving \(C\).
If \(\dim_{\F_4}W=2\), the quotient has dimension two.
A zero action could not generate \(W\) from its single alternating
coefficient. For a nonzero square-zero action \(P_1\), its commuting
square-zero partners are scalar multiples of it. Change quotient basis
so \(P_2=0\), and normalize
\(a_{12}=z,P_1z=y,P_1y=0\). This gives \(N\).
The displayed intrinsic dimensions distinguish the three classes.
\end{proof}

\subsection{Central descent with the smallest nontrivial abelianization}

We now assume solvability and fix \(\dim_{\F_2}(V/W)=4\).
This fixes the smallest abelianization dimension of a nonabelian
nilpotent polar algebra. The following theorem closes the stratum in
all ambient dimensions.

\begin{theorem}[A closed seven-class chain]\label{thm:seven-chain}
Among solvable polar solutions with binary abelianization dimension four,
there is exactly one conjugacy class in each of dimensions
\(4,6,8,10,12,14,16\), and no class in any other dimension.
Their polar classes are \(1,2,3,3,4,5,5\).
The extension-space ranks in the proof are exact finite computations,
certified as described in Appendix~\ref{app:verification}.
\end{theorem}

The representatives are the successive prefixes of the following map,
with coordinates ordered as \((a,b,c,d,e,h,r,s)\):
\begin{equation}\label{eq:seven-coordinates}
\begin{aligned}
 a'&=\omega a, &b'&=\omega b,\\
 c'&=\omega c+ab,\\
 d'&=\omega d+ac+a^2b,\\
 e'&=\omega e+bc+ab^2,\\
 h'&=\omega h+ae+bd,\\
 r'&=\omega r+A,& A&=ac^2+a^2e+ah+cd,\\
 s'&=\omega s+B,& B&=bc^2+b^2d+bh+ce.
\end{aligned}
\end{equation}
Each prefix is binary-quadratic. Every new cocycle \(q\) satisfies
\(q(g(x))=\omega^2q(x)\) for the preceding prefix \(g\), proving
\eqref{eq:functional} recursively. For the last two expressions this
follows by substitution; the extra terms cancel in pairs or have
coefficient \(1+\omega+\omega^2=0\).

\begin{table}[ht]
\centering
\caption{Intrinsic invariants of the seven prefixes. All dimensions are binary.}
\label{tab:chain-invariants}
\begin{tabular}{@{}rrrrl@{}}
\toprule
\(n\)&\(\dim W\)&\(\dim Z\)&Class&Lower central dimensions\\
\midrule
4&0&4&1&\((4,0)\)\\
6&2&2&2&\((6,2,0)\)\\
8&4&2&3&\((8,4,2,0)\)\\
10&6&4&3&\((10,6,4,0)\)\\
12&8&2&4&\((12,8,6,2,0)\)\\
14&10&2&5&\((14,10,8,4,2,0)\)\\
16&12&4&5&\((16,12,10,6,4,0)\)\\
\bottomrule
\end{tabular}
\end{table}

\begin{proof}[Proof of Theorem~\ref{thm:seven-chain}]
The proof proceeds from descent to extensions and then to closure.
Every member descends to the abelian four-dimensional solution.
At each prefix we compute the complete quadratic extension quotient;
kernel scalings and quotient automorphisms put all its nonzero vectors
into a single orbit. This gives a unique next class until the final
extension quotient vanishes. Descent then turns that last calculation
into a bound on every ambient dimension.

\emph{Descent.}
Solvability implies nilpotency. In a nonabelian member of this stratum,
\(Z(L)\subseteq W\). Otherwise the linear scalar action on \(Z(L)\)
would give a central invariant \(\F_4\)-line disjoint from \(W\).
Its quotient would be a nonabelian nilpotent polar algebra with
binary abelianization dimension two, contradicting
Lemma~\ref{lem:two-generators}.
Quotient by any central \(\F_4\)-line in \(W\). The dimension drops
by two, the abelianization dimension stays four, and a quadratic
solution is inherited. Repetition ends at the unique abelian
four-dimensional solution. Every member is therefore reached by
successive central \(\F_4\)-extensions.

\emph{The complete extension space.}
For a fixed quotient solution \(g\), every such extension has the form
\[
 F(x,z)=(g(x),\omega z+q(x)),\qquad
 q(g(x))=\omega^2q(x),\quad q(0)=0,
\]
where \(q\) is binary-quadratic.
Let \(E_g\) be this \(\F_4\)-vector space and define
\begin{equation}\label{eq:extension-quotient}
 H_g=E_g/\delta\Hom_{\F_2}(V,\F_4),\qquad
 \delta(t)=t\circ g+\omega t.
\end{equation}
This quotient accounts for every binary-linear change of splitting.
Kernel scalings and quotient automorphisms act by scaling and pullback.
The zero class splits off a central scalar line and increases the
binary abelianization dimension to six.
Conversely, if an extension has its new line disjoint from its derived
ideal, choose an \(\F_4\)-complement to that line in the abelianization
and take its inverse image. It is an invariant complement containing
the derived ideal, so the extension splits and its class in \(H_g\)
is zero. Thus exactly the nonzero classes remain in our stratum.

\emph{Ranks and representatives.}
The complete results are in Table~\ref{tab:extension-ranks}.
The equation for \(q\) has binary degree at most four, so values at
inputs of binary weight at most four recover every coefficient by
Boolean M\"obius inversion. The certificates verify the ranks and
the spanning of each displayed quotient. The splitting-image dimension
also follows directly: \(\delta(t)=0\) precisely for the linear
maps \(t\) that annihilate \(W\) and intertwine the intrinsic
\(\omega\)-actions. This kernel has binary dimension four, so
\(\dim\im\delta=2n-4\).

\begin{table}[ht]
\centering
\caption{Exact central-extension spaces. Here \(n\) is the dimension of \(g\).}
\label{tab:extension-ranks}
\begin{tabular}{@{}rrrrrrl@{}}
\toprule
\(n\)&Unknowns&Rank&\(\dim_{\F_2}E_g\)&\(\dim_{\F_2}\im\delta\)&
\(\dim_{\F_4}H_g\)&Basis of \(H_g\)\\
\midrule
4&20&14&6&4&1&\(ab\)\\
6&42&30&12&8&2&\(ac+a^2b,\ bc+ab^2\)\\
8&72&58&14&12&1&\(bc+ab^2\)\\
10&110&92&18&16&1&\(ae+bd\)\\
12&156&132&24&20&2&\(A,B\)\\
14&210&184&26&24&1&\(B\)\\
16&272&244&28&28&0&--\\
\bottomrule
\end{tabular}
\end{table}

\emph{Projective orbits.}
When \(H_g\) is one-dimensional, kernel scalings identify all its
nonzero classes. There are two two-dimensional cases.
For \(n=6\), put \(u=ac+a^2b,v=bc+ab^2\).
The following automorphisms of \(g\) act on their classes by a swap,
a nontrivial diagonal scaling, and a shear:
\begin{align*}
 &(a,b,c)\mapsto(b,a,c),\\
 &(a,b,c)\mapsto(\lambda a,b,\lambda c),
       \qquad\lambda\in\F_4^*,\\
 &(a,b,c)\mapsto(a+b,b,c+b^2).
\end{align*}
Their actions are \((u,v)\mapsto(v,u)\),
\((\lambda^2u,\lambda v)\), and \((u+v,v)\).
They are transitive on the five projective points of \(\F_4^2\).

For \(n=12\), the corresponding automorphisms are
\begin{align*}
 &(a,b,c,d,e,h)\mapsto(b,a,c,e,d,h),\\
 &(a,b,c,d,e,h)\mapsto
       (\lambda a,b,\lambda c,\lambda^2d,\lambda e,\lambda^2h),\\
 &(a,b,c,d,e,h)\mapsto(a+b,b,c+b^2,d+e,e,h).
\end{align*}
They send \((A,B)\) to \((B,A)\), \((A,\lambda^2B)\), and
\((A+B+ab+b^2,B+b^2)\). The extra terms are splitting changes:
\(ab=\delta(c)\), \(b^2=\delta(b^2)\).
The projective action is again transitive.

\emph{Closure.}
Starting with the abelian four-dimensional map, every nonzero extension
class therefore yields the unique next prefix, through dimension sixteen.
At that point \(H_g=0\), so no next member exists. Descent rules out
every higher dimension as well.
The coordinate brackets give Table~\ref{tab:chain-invariants}.
Their second derived dimensions in dimensions fourteen and sixteen are
two and four; these are the two nonmetabelian members.
\end{proof}

The six-dimensional prefix is the central example already present in
\cite[Lemma~10.3(1)]{Zhang1998}. General central recursion also appears
there. The quadratic eigenspaces, linear-conjugacy orbits, and the
vanishing of the final extension quotient are the ingredients that
close the present seven-class range.

\section{Further structural questions}
\label{sec:questions}

The structure theory suggests three directions, each with a concrete
starting point and a different obstruction to overcome.

\paragraph{Can a nonsolvable polar algebra occur?}
Corollary~\ref{cor:simple-reduction} reduces this to a simple polar
example. Such an algebra must have even binary dimension, satisfy
\eqref{eq:four-variable}, and admit a commuting quadratic refinement.
The matrix and Witt obstructions give two tests for candidate algebras.
A general obstruction would settle the question negatively; a simple
polar example would produce solutions beyond the nilpotent regime.
The main issue is to construct or exclude a quadratic refinement
satisfying the full functional equation.

\paragraph{What is the dimension cost beyond the metabelian case?}
The fourteen-dimensional member of
Theorem~\ref{thm:seven-chain} already has class five, whereas the
metabelian minimum at that class is thirty-six. Thus the all-solution
minimum at class five is at most fourteen, and lower bounds for that
problem must account for nonabelian derived ideals. The closed
abelianization-four stratum supplies exact reference cases. A general
theory would need to control further extension patterns as the
abelianization grows, beyond the commuting module actions used here.

\paragraph{Which structures attain the sharp metabelian bounds?}
Theorem~\ref{thm:extremizers} gives at least \(c-2\) compatible equality
cases in class \(c\ge4\). Classifying all equality cases requires
determining which module quotients and cocycles attain the bounds,
and when their resulting maps are conjugate. The contraction quotients
in \eqref{eq:terminal-quotient}, the filtered action spaces
\(\mathcal A_k\), and the morphism correspondence provide tools for
this problem and for computing the extremizers' automorphism groups.

There is also a geometric interpretation of every result: nonzero
orbits define an oriented binary line spread, and linear conjugacy
preserves the underlying additive geometry. The polar Lie algebra
records degree-two interactions between those oriented lines.
Its center has a direct functional meaning. For \(a\in Z(L)\),
\[
 f(x+a)+f(x)=f(a)\qquad\text{for every }x.
\]
Consequently nilpotent models have constant differences in nonzero
central directions. Their structure is particularly suited to the
algebraic and geometric classification questions developed here.

\clearpage
\appendix
\section{Exact certificates for the central-extension ranks}
\label{app:verification}

This appendix specifies the finite rank proof in
Theorem~\ref{thm:seven-chain}. The general correspondence and both
all-class extremal bounds are proved independently in the main text.

\subsection{Encoding and equations}

Use the binary basis \(1,\omega\) of \(\F_4\), with
\(\omega^2=\omega+1\). For an \(n\)-dimensional binary domain,
the nonconstant ANF monomials of degree at most two are
\(x_i\) and \(x_ix_j\), \(i<j\).
A candidate \(\F_4\)-valued \(q\) therefore has
\(2(n+\binom n2)=n(n+1)\) binary unknowns.
The map \(g\) is stored by the binary output vector of each such
monomial, with the scalar-coordinate order of
\eqref{eq:seven-coordinates}.

For each binary point \(x\), evaluating
\(q(g(x))+\omega^2q(x)\) gives two binary rows in these unknowns.
Their degree is at most four. Requiring all rows at inputs of weights
one through four is equivalent to requiring the entire function to
vanish: the value at zero is already zero, and triangular Boolean
M\"obius inversion recovers every coefficient of degree at most four.

\subsection{What each certificate proves}

For each \(n\in\{4,6,8,10,12,14,16\}\), the supplied certificate contains:
\begin{enumerate}
\item the quadratic coefficients of the quotient map and the monomial order;
\item a stated number of independent rows, each with its originating
      input point and output bit;
\item a kernel basis annihilated by every weight-at-most-four row;
\item the actual vectors \(\delta(t)\) for a basis of binary-linear
      changes of splitting;
\item the named representatives of \(H_g\) in
      Table~\ref{tab:extension-ranks}.
\end{enumerate}
The independent checker reconstructs the rows from the stored map,
checks row independence over \(\F_2\), checks the entire kernel on all
required inputs, and verifies that row rank plus kernel dimension is
the number of unknowns. It then reconstructs the splitting-change
vectors and verifies that their span together with the stated quotient
representatives is exactly the kernel.
Thus both the upper and lower rank bounds are certified.

These are exact, reproducible algebraic checks; they are not
proof-assistant formalizations.

\subsection{Files and replay}

The \path{verification} directory contains the seven certificates,
their SHA-256 manifest, and three checkers. From the source-project
directory, run:
\begin{verbatim}
cd verification
python check_chain_rank_certificates.py
python audit_metabelian_symbolic_interfaces.py
python manuscript_transcription_checks.py
\end{verbatim}
Use Python 3 and, for the first program, SymPy. The rank checker
authenticates each record, rebuilds the equations from the stored map,
and uses exact \(\F_2\) linear algebra in SymPy's domain-matrix
implementation; it imports no generator code. The other programs use
the standard library: the second checks scalar weights, binary degree,
and morphism interfaces; the third checks coordinate transcription,
intrinsic dimensions, and projective actions. Each prints JSON results
with status \texttt{PASS} and writes reports under \path{verification}.
Expected ranks and quotient dimensions are in
Table~\ref{tab:extension-ranks}.

\clearpage
\begin{thebibliography}{99}

\bibitem{Drensky2021}
V.~Drensky,
Weak polynomial identities and their applications,
\emph{Communications in Mathematics} \textbf{29} (2021), 291--324.
\href{https://doi.org/10.2478/cm-2021-0022}{doi:10.2478/cm-2021-0022}.

\bibitem{KovacsStohr2013}
L.~G.~Kov\'acs and R.~St\"ohr,
\emph{Free centre-by-metabelian Lie algebras in characteristic 2},
MIMS EPrint 2013.19, author manuscript, 2013.
\url{https://eprints.maths.manchester.ac.uk/1971/1/KS_2012_submit.pdf}.

\bibitem{LieAlgDB}
LieAlgDB package documentation,
\emph{A description of the Lie algebras that are contained in the package},
Chapter~3, Section~3.6.
\url{https://gap-packages.github.io/liealgdb/doc/chap3.html}.

\bibitem{Li2026}
B.~Li,
\emph{Norm-fiber densities and relative traces in characteristic two},
preprint with supplementary material, 2026.
\href{https://doi.org/10.5281/zenodo.23757428}{doi:10.5281/zenodo.23757428}.

\bibitem{Liu2019}
X.~Liu, Y.~Chen, Y.~Xu, and Z.~Sun,
Triple-cycle permutations over finite fields of characteristic two,
\emph{International Journal of Foundations of Computer Science}
\textbf{30} (2019), no.~2, 275--292.
\href{https://doi.org/10.1142/S0129054119500059}{doi:10.1142/S0129054119500059}.

\bibitem{LopatkinZusmanovich2019}
V.~Lopatkin and P.~Zusmanovich,
\emph{Commutative Lie algebras and commutative cohomology in characteristic 2},
preprint, 2019.
\url{https://arxiv.org/abs/1907.03690}.

\bibitem{Lu2023}
W.~Lu, X.~Wu, Y.~Wang, and X.~Cao,
A general construction of regular complete permutation polynomials,
\emph{Designs, Codes and Cryptography} \textbf{91} (2023), 2627--2647.
\href{https://doi.org/10.1007/s10623-023-01224-6}{doi:10.1007/s10623-023-01224-6}.

\bibitem{Pakianathan2007}
J.~Pakianathan and E.~Yal\c{c}{\i}n,
Quadratic maps and Bockstein closed group extensions,
\emph{Transactions of the American Mathematical Society}
\textbf{359} (2007), 6079--6110.
\href{https://doi.org/10.1090/S0002-9947-07-04346-2}{doi:10.1090/S0002-9947-07-04346-2}.

\bibitem{PapadimaSuciu2004}
\c{S}.~Papadima and A.~I.~Suciu,
Chen Lie algebras,
\emph{International Mathematics Research Notices}
\textbf{2004} (2004), no.~21, 1057--1086.
\href{https://doi.org/10.1155/S1073792804132017}{doi:10.1155/S1073792804132017}.

\bibitem{Samardjiska2017}
S.~Samardjiska and D.~Gligoroski,
Quadratic permutation polynomials, complete mappings and mutually orthogonal
Latin squares,
\emph{Mathematica Slovaca} \textbf{67} (2017), no.~5, 1129--1146.
\href{https://doi.org/10.1515/ms-2017-0037}{doi:10.1515/ms-2017-0037}.

\bibitem{Strade2006}
H.~Strade,
\emph{Lie algebras of small dimension}, preprint, 2006.
\url{https://arxiv.org/abs/math/0601413}.

\bibitem{VaughanLee2006}
M.~Vaughan-Lee,
Simple Lie algebras of low dimension over GF(2),
\emph{LMS Journal of Computation and Mathematics}
\textbf{9} (2006), 174--192.
\href{https://doi.org/10.1112/S1461157000001248}{doi:10.1112/S1461157000001248}.

\bibitem{Xu2022}
X.~Xu, X.~Zeng, and S.~Zhang,
Regular complete permutation polynomials over \(\F_{2^n}\),
\emph{Designs, Codes and Cryptography} \textbf{90} (2022), 545--575.
\href{https://doi.org/10.1007/s10623-021-00992-3}{doi:10.1007/s10623-021-00992-3}.

\bibitem{Zhang1998}
X.~Zhang,
\emph{Nilpotent Algebras with Maximal Class in Congruence Modular Varieties},
Ph.D. thesis, University of Manitoba, 1998.
\url{https://mspace.lib.umanitoba.ca/bitstream/handle/1993/1662/NQ32041.pdf}.

\end{thebibliography}

\end{document}
